% \section{Abstract}

\begin{abstract}


While programmers know that memory representation of data structures can have significant effects
on performance, compiler support to {\em optimize} the layout of those structures is an under-explored field.
Prior work has optimized the layout of individual, \textit{non-recursive} structures without considering how
collections of those objects in linked or \textit{recursive} data structures are laid out.
%The internal layout of constituent objects in linked data structures is an under-explored field.

%or focused on arranging the 
%placement of linked data structures without considering the internal layout of the constituent objects.

This work introduces \system, a compiler that optimizes the layouts of algebraic datatypes, with a special focus
on producing highly optimized, {\em packed} data layouts where recursive structures can be traversed with minimal
pointer chasing. \system performs an analysis of how a recursive ADT is used across functions to choose a \textit{global} layout
that promotes simple, strided access for that ADT in memory. It does so by building and solving a constraint system
to minimize an abstract cost model, yielding a predicted efficient layout for the ADT. \system then builds on top
of \gibbon, a prior compiler for packed, mostly-serial representations, to synthesize optimized ADTs. We show experimentally 
that \system is able to choose optimal layouts across a series of microbenchmarks and case studies, outperforming both 
\gibbon's baseline approach, as well as \mlton, a Standard ML compiler that uses traditional
pointer-heavy representations.


%\bf{ABSTRACT 2}
%
%Compilers treat the layout of fields in memory as a matter of
%preestablished convention, not as a free variable to optimize for a
%particular input program.
%%
%This rigidity incurs an opportunity cost, especially as we move to
%consider more compact in-memory representations for data; for example,
%as found in the \gibbon compiler which uses compact, mostly-serialized
%heap formats.
%
%This paper introduces \system, a compilation approach for functional
%programs which uses a solver to determine the optimal-cost ordering of
%data fields, optimizing for a particular traversal function over that
%data.
%%
%An experimental evaluation and case study shows that this approach
%yields an \auditme{XYZ} speedup over the compact-data approach alone,
%and an \auditme{XYZ} speedup over the mature GHC compiler with its
%traditional, pointer-heavy, non-compact data representations.
  
\end{abstract}
